\chapter{Writing a four-hour exam and the proof portfolio}
\index{exam technique}\index{proof portfolio}

The previous sixteen chapters taught the mathematics; this one teaches
the interface where mathematics meets a clock and a grader. A contest
score is not a sum over problems you solved --- it is a sum over proofs
the grader could \emph{follow}. After this chapter you can budget four
hours, write a proof a grader marks, and maintain a twenty-proof
portfolio that measures itself.

\section{What a timed contest tests}\index{contest!what it tests}
Homework tests whether you can solve a problem eventually. A four-hour
contest tests whether you can solve \emph{enough, fast, and legibly}.
The skill that homework never exercises is abandonment: knowing that a
problem you \emph{could} solve in three hours is a loss if the other
five were solvable in thirty minutes each.

\begin{example}[Triage by rate, not by difficulty]
An exam offers problems worth 7 points each: $P_1$ looks solvable in
40 minutes, $P_2$ in 90, $P_3$ is interesting and you have spent 25
minutes already with no progress. Rank the next hour.
\par\smallskip
\textbf{Strategy.} Score is additive; time is the constraint ---
maximize expected points per \emph{hour}, not per problem, and treat
sunk cost as zero.
\par\smallskip
\begin{solution}
$P_1$: $\approx 7/\tfrac{2}{3} = 10.5$ pts/hr; $P_2$:
$\approx 7/1.5 \approx 4.7$; $P_3$: unknown remainder, likely $>60$
min more for at most 4.7. The 25 minutes on $P_3$ are gone either way
--- they do not raise its rate. Order: $P_1$ first, then re-evaluate
$P_2$ vs returning to $P_3$ with fresh eyes (Figure~\ref{fig:triage}).
\end{solution}
\end{example}

\begin{figure}[ht]\centering
\begin{tikzpicture}[x=1.05cm,y=0.6cm]
  % points-per-hour bars for a 3-problem triage decision
  \foreach \x/\h/\l in {0.5/4.5/{P1: doable, 40 min},
                        2.7/3.2/{P2: risky, 90 min},
                        4.9/1.6/{P3: sunk-cost trap}}{
    \draw[fill=accent!30] (\x-0.6,0) rectangle (\x+0.6,\h);
    \node[font=\small] at (\x,-0.7) {\l};
  }
  \draw[->] (-0.6,0) -- (6.2,0);
  \draw[->] (-0.6,0) -- (-0.6,5.2) node[above,font=\small] {est.\ pts/hr};
  \node[font=\small,align=left,anchor=west] at (6.4,3) {order by\emph{rate},\\not difficulty};
\end{tikzpicture}

\caption{Triage orders by expected points per hour, not by how close a
problem \emph{feels}. Notice $P_3$'s sunk cost buys nothing --- only
the remaining estimate counts.}
\label{fig:triage}
\end{figure}

\section{The four-hour budget}\index{contest!budget}
Triage first, always: read every problem before writing any
(Figure~\ref{fig:gantt}). First pass solves the problems you already
know how to finish --- typically half the marks for a quarter of the
time. Deep work then goes to \emph{one} problem at a time, chosen by
point value times your probability of finishing, not by how interesting
it looks. A running clock answers ``should I abandon this?'' ---
fifteen minutes without a new idea is a new-idea problem, not a
persistence problem.
\begin{figure}[ht]\centering
% Figure: time-budget strip for a 4-hour olympiad.
% Source for Figure~\ref{fig:gantt}. TikZ only; regenerable by compiling main.tex.
\begin{tikzpicture}[
  lab/.style={font=\footnotesize\sffamily,text=slate,align=center},
  seg/.style={draw=ink,minimum height=0.42cm,inner sep=1.5pt,font=\scriptsize\sffamily,align=center},
]
\node[seg,fill=rulegray!40,minimum width=1.0cm] at (0.5,0) {read all\\10 min};
\node[seg,fill=accent!25,minimum width=2.6cm] at (2.3,0) {first pass: P1--P2 owned\\80 min};
\node[seg,fill=accent!45,minimum width=2.4cm] at (4.8,0) {deep work: one problem\\fully};
\node[seg,fill=rulegray!60,minimum width=1.4cm] at (6.7,0) {write-ups\\+ check\\30 min};
\draw[rulegray,->] (0,-0.5) -- (7.6,-0.5);
\foreach \x/\t in {0/0, 1/30, 3/90, 5/150, 6/210, 7.5/240}
  \node[lab] at (\x,-0.75) {\t};
\node[lab] at (3.8,-1.25) {minutes elapsed --- the last 30 are never ``extra work''};
\end{tikzpicture}

\caption{A four-hour budget. Notice the protected last thirty minutes:
checking is scheduled, not squeezed.}
\label{fig:gantt}
\end{figure}
\begin{example}[Two budgets compared]\index{budget comparison}
$15$ answer-only problems in $3$ hours vs $6$ proofs in $4.5$ hours:
how should the allocation differ?

\textbf{Strategy.} Budget by expected value: (points) $\times$
(probability of completion), re-evaluated at each checkpoint.

\textbf{Solution.} Answer-only: uniform-ish spread, abandon at the
first stall ($15 \times 12$ min each; a $20$-minute problem is a losing
one). Proof contest: front-load certainty --- the safest proof goes
first, the 50/50 problem gets a hard time cap. The budgets differ in
\emph{shape}: the first is flat, the second is a waterfall.
\end{example}

\section{Write-up standards}\index{write-up standards}
A grader reads top-to-bottom at speed. So: \textbf{claim first} ---
state what you will prove before proving it. \textbf{Name the lemmas}
--- a lemma with a name is a checkpoint the grader can tick.
\textbf{Display the math} --- buried algebra is invisible algebra.
\textbf{Draw the diagram} --- one figure replaces a paragraph.
\textbf{Follow checking conventions} --- the argument order is
invention order reversed (Figure~\ref{fig:pipeline}).
\begin{figure}[ht]\centering
% Figure: strategy -> skeleton -> details write-up pipeline.
% Source for Figure~\ref{fig:pipeline}. TikZ only; regenerable by compiling main.tex.
\begin{tikzpicture}[
  lab/.style={font=\footnotesize\sffamily,text=slate,align=center},
  box/.style={draw=ink,rounded corners=2pt,minimum height=0.9cm,font=\footnotesize\sffamily,align=center,inner sep=4pt},
]
\node[box,fill=rulegray!30] (a) at (0,0) {strategy\\(scratch paper:\\the idea)};
\node[box,fill=accent!25] (b) at (3.1,0) {skeleton\\(claim, lemmas,\\order of steps)};
\node[box,fill=accent!45] (c) at (6.2,0) {details\\(full sentences,\\displayed math)};
\draw[rulegray,thick,->] (a) -- (b);
\draw[rulegray,thick,->] (b) -- (c);
\node[lab,align=center]{the artifact is built in three passes ---\\
never submitted in discovery order};
\end{tikzpicture}

\caption{The write-up pipeline. Notice what changes at each pass: the
idea stays fixed while the \emph{reader's needs} are added.}
\label{fig:pipeline}
\end{figure}
\begin{example}[Scratch vs submission]\index{write-up!example}
Chapter 16's difference-reduction problem, shown as (a) the scratch
note a solver actually writes and (b) the text a grader receives.

\textbf{Strategy.} Same content, different artifact: scratch says
``try small cases, gcd stays!''; the submission turns that discovery
into claim--lemma--proof order.

\textbf{Solution.} \emph{Scratch:} ``$12,18,30,42\to$ differences,
gcd seems to survive, sum shrinks so it stops --- write it up.''
\emph{Submission:} \textbf{Claim.} The survivor equals
$\gcd(a_1,\dots,a_n)$. \textbf{Lemma (Termination).} The sum strictly
decreases, so the process ends. \textbf{Lemma (Invariant).}
$\gcd$ is preserved since $\gcd(a-b,b)=\gcd(a,b)$. \textbf{Proof.}
Termination gives a final state; the invariant evaluates it.\ \qed\
The scratch line contained the whole solution; the submission made it
gradeable.
\end{example}

\section{Partial credit and fake-solving}\index{partial credit}
Graders reward \emph{progress they can verify}: a stated claim, a
proved lemma, a correctly argued case. They cannot reward ``I was
close''. \emph{Fake-solving} --- asserting the hard step as ``clearly''
or ``similarly'' --- earns nothing and costs the time you spent
writing it (Figure~\ref{fig:sidebyside}). The honest partial write-up
that proves the easy case and names the gap out-scores a confident
bluff every time.
\begin{figure}[ht]\centering
% Figure: bad/good write-up side by side (our own marked-up text).
% Source for Figure~\ref{fig:sidebyside}. TikZ only; regenerable by compiling main.tex.
\begin{tikzpicture}[
  lab/.style={font=\footnotesize\sffamily,text=slate,align=center},
  pg/.style={draw=ink,minimum width=3.2cm,minimum height=2.4cm,align=left,font=\scriptsize,inner sep=4pt,text width=3.0cm},
]
\node[pg,fill=red!5] (b) at (0,0) {``Clearly the answer is 6.\\
Trying 3 cases works, so\\
we are done. Also the\\
pattern continues.''\\
\textcolor{red}{no claim, no proof,\\cases asserted not argued}};
\node[pg,fill=green!6] (g) at (3.9,0) {\textbf{Claim.} The survivor is\\
$\gcd(a_1,\dots,a_n)$.\\
\textbf{Proof.} Sum drops each\\
move \ldots\ gcd is preserved\\
since $\gcd(a-b,b)=\gcd(a,b)$.\\
\qed};
\node[lab] at (1.95,-1.7) {same mathematics --- only one earns marks:\\
claim first, named steps, argued cases};
\end{tikzpicture}

\caption{Bad vs good write-up of the same mathematics. Notice the bad
one isn't shorter --- it is \emph{unverifiable}, which reads as zero.}
\label{fig:sidebyside}
\end{figure}
\begin{example}[Autopsy of a bad write-up]\index{fake-solve}
``Clearly the survivor is the gcd. Trying a few cases shows it works,
so we are done.''

\textbf{Strategy.} Grade it: mark what a grader could verify.

\textbf{Solution.} No claim environment, no termination argument, the
gcd-preservation step asserted not shown, and ``a few cases'' carries
zero weight against an infinite family. Verifiable content: none ---
0/7 (on the 0--7 USAMO scale). The fix is not more words; it is the two named lemmas.
\end{example}

\section{Checking your work}\index{checking}
Last thirty minutes, in order of value: (1) re-read each submission's
claim --- does the proof actually end at it? (2) check boundary cases
--- the $n=0$, empty-set, single-element input. (3) verify find-all
answers in both directions: candidates found \emph{and} verified.
(4) recompute any arithmetic that appeared only once on the page
(Figure~\ref{fig:checkloop}).

\begin{figure}[ht]\centering
\begin{tikzpicture}[x=1cm,y=1cm,>=stealth,every node/.style={font=\small,draw,rounded corners=3pt,minimum height=0.9cm}]
  \node[fill=accent!20] (solve) at (0,0) {solve};
  \node[fill=graybg] (dim) at (2.9,0.9) {dimensions/units?};
  \node[fill=graybg] (edge) at (2.9,-0.9) {edge cases?};
  \node[fill=graybg] (plug) at (5.8,0) {plug back in};
  \node[fill=good!25] (submit) at (8.8,0) {submit};
  \draw[->] (solve) -- (dim);
  \draw[->] (solve) -- (edge);
  \draw[->] (dim) -- (plug);
  \draw[->] (edge) -- (plug);
  \draw[->] (plug) -- (submit);
  \draw[->,dashed] (plug) to[bend left=45] node[below,font=\small,draw=none]{any check fails} (solve);
\end{tikzpicture}

\caption{The check loop: three cheap checks before submit, and any
failure routes back to solving. Notice checks are ordered by cost ---
dimensions are seconds, plug-back is the thorough audit.}
\label{fig:checkloop}
\end{figure}

\begin{example}[A check that catches the answer]
A solution claims the sequence $a_{n+1}=2a_n+1$, $a_0=0$ has closed
form $a_n = 2^{n-1}-1$. Run the check loop.
\par\smallskip
\textbf{Strategy.} Cheap check first: the boundary value $n=1$ must
already satisfy it.
\par\smallskip
\begin{solution}
$a_1 = 1$ from the recurrence; the formula gives
$2^0-1=0$ --- contradiction at the first check, costing seconds.
The correct form is $a_n = 2^n - 1$: $2^1-1=1$, $2^2-1=3$ match the
recurrence's $a_1, a_2$, and induction confirms. One boundary case
turned a confidently-wrong answer into a fix before submission.
\end{solution}
\end{example}

\section{The 20-proof portfolio}\index{proof portfolio}
Twenty proofs, written and self-graded on the rubric
(Figure~\ref{fig:rubric}): claim stated, lemmas named, cases argued,
math displayed, self-check passed --- each scored 0/1/2. Keep the
lowest-scoring proof, rewrite it, and only then add a twenty-first.
The portfolio grows by \emph{replacing} your weakest work, not by
accumulating new work.
\begin{figure}[ht]\centering
% Figure: portfolio rubric table as a diagram.
% Source for Figure~\ref{fig:rubric}. TikZ only; regenerable by compiling main.tex.
\begin{tikzpicture}[
  lab/.style={font=\footnotesize\sffamily,text=slate,align=center},
  hd/.style={draw=ink,fill=slate!15,font=\scriptsize\sffamily\bfseries,minimum width=2.2cm,minimum height=0.34cm,inner sep=2pt},
  cl/.style={draw=ink,font=\scriptsize,minimum width=2.2cm,minimum height=0.34cm,inner sep=2pt},
]
\node[hd] at (-2.2,0.6) {rubric item};
\node[cl] at (-2.2,0.25) {claim stated first};
\node[cl] at (-2.2,-0.1) {lemmas named};
\node[cl] at (-2.2,-0.45) {every case argued};
\node[cl] at (-2.2,-0.8) {displayed math};
\node[cl] at (-2.2,-1.15) {self-check pass};
\node[hd] at (0.2,0.6) {0 / 1 / 2};
\node[cl] at (0.2,0.25) {absent/buried/first};
\node[cl] at (0.2,-0.1) {none/some/all};
\node[cl] at (0.2,-0.45) {asserted/mixed/argued};
\node[cl] at (0.2,-0.8) {inline/mixed/displayed};
\node[cl] at (0.2,-1.15) {skipped/errors/clean};
\node[lab,align=center]{score each of 20 proofs; rewrite the\\
lowest-scoring dimension until it's a 2};
\end{tikzpicture}

\caption{The portfolio rubric. Notice it is mechanical on purpose ---
a rubric that needs judgment can't grade itself.}
\label{fig:rubric}
\end{figure}

\section{Code: a linter, not a grader}\index{portfolio\_check}
\texttt{portfolio\_check.py} scans your proof files for the mechanical
parts of the rubric --- a claim environment, named lemmas, a proof
environment, a \verb|\qed|. It cannot see fake-solves; it only sees
missing structure. That is the honest division of labor.
\begin{checkpoint}
Predict which sample file passes. The good one carries a theorem, a
named lemma, a proof and \verb|\qed|; the bad one is prose. Then run
the linter.
\end{checkpoint}
\begin{lstlisting}
sample_proofs/bad_proof.tex:
   missing: claim first
   missing: named lemma
   missing: proof env
   missing: qed
sample_proofs/good_proof.tex: OK
\end{lstlisting}
\texttt{test\_17\_portfolio.sh} asserts both lines.

\section{Summary}
\begin{center}\small
\begin{tabular}{@{}ll@{}}
\toprule
\textbf{Minute} & \textbf{Action} \\
\midrule
0--10 & read everything; rank by value$\times$probability \\
10--90 & first pass: the problems you own \\
90--210 & deep work, one problem at a time, hard caps \\
210--240 & write-up check: claim, boundary, both directions \\
after & portfolio: rewrite the weakest proof \\
\bottomrule
\end{tabular}
\end{center}

\section{Common mistakes}\index{exam technique!mistakes}
\begin{itemize}
\item \emph{Discovery-order writing.} The order you found the proof is
almost never the order a reader needs. Reorder.
\item \emph{Half a find-all.} Finding the candidates is half the
answer; the other half is verifying each one solves it.
\item \emph{Polishing problem 1 while 4--6 are blank.} A perfect early
proof plus three empty pages loses to three rough proofs. The budget
is the strategy.
\end{itemize}

\section*{Exercises}
\subsection*{Inline}
\begin{exercise}[easy]\label{exr:ex1}
A 3-problem contest gives 90 minutes. State a budget for a student who
expects to solve two and struggle on one.
\end{exercise}
\begin{solution}
5-min read, $2\times30$ min on the two owned problems, $15$ min capped
on the third for partial credit, $10$ min checking.
\end{solution}
\begin{exercise}[medium]\label{exr:ex2}
Rewrite this ``proof'' to passing standard: ``$n^2$ is even for even
$n$, so $n^2+n$ is even.''
\end{exercise}
\begin{solution}
Claim first; split parity as a named lemma covering \emph{both} cases;
close with the sum argument and \qed.
\end{solution}
\begin{exercise}[medium]\label{exr:ex3}
Your write-up solved a find-all equation by listing three solutions
and noting they work. What is missing and how many rubric points does
it cost?
\end{exercise}
\begin{solution}
The ``no other solutions'' direction --- a completeness proof. On the
rubric the ``every case argued'' item scores 0 of its 2 points, so the
omission costs 2 of 10.
\end{solution}
\subsection*{Additional}
\begin{exercise}[easy]\label{exr:ea1}
List the three checks to run in the last 30 minutes, highest value
first.
\end{exercise}
\begin{solution}
Claim-matches-proof; boundary cases; both directions on find-alls.
\end{solution}
\begin{exercise}[easy]\label{exr:ea2}
Why does a rubric need to be mechanical?
\end{exercise}
\begin{solution}
So self-assessment is repeatable --- a rubric requiring judgment
re-imports the bias it exists to remove.
\end{solution}
\begin{exercise}[easy]\label{exr:ea3}
In 15 answer-problems/3h, what does ``abandon'' look like concretely?
\end{exercise}
\begin{solution}
A hard cap per problem ($\approx$12 min); at the cap with no new idea,
write the best guess and move.
\end{solution}
\begin{exercise}[medium]\label{exr:ea4}
Timed rewrite: take your Chapter 16 difference-reduction solution and
produce a submission draft in 20 minutes; score it on the rubric.
\end{exercise}
\begin{solution}
Expect claim+two lemmas+proof; a missing termination lemma is the
usual 1-point leak.
\end{solution}
\begin{exercise}[medium]\label{exr:ea5}
Explain why fake-solving loses \emph{expected} marks even when the
fake step is true.
\end{exercise}
\begin{solution}
Asserted steps are unverifiable (0 credit), and the bluff hides the
real partial credit available for the argued parts.
\end{solution}
\begin{exercise}[medium]\label{exr:ea6}
Design the portfolio update rule and justify ``replace weakest'' over
``add newest''.
\end{exercise}
\begin{solution}
A fixed-size portfolio forces the floor up; adding only lets weak work
persist.
\end{solution}
\begin{exercise}[medium]\label{exr:ea7}
Implement-or-derive: extend \texttt{portfolio\_check.py} to also flag
``clearly'', ``obviously'', and ``similarly'' as fake-solve markers.
\end{exercise}
\begin{solution}
Add a fourth regex group over those words; each hit is a 0-point
marker in the report line.
\end{solution}
\begin{exercise}[hard]\label{exr:ea8}
Timed: rewrite any earlier chapter's worked example as a full
submission in 25 minutes, self-score, then rewrite once more to raise
the lowest dimension.
\end{exercise}
\begin{solution}
The second pass should visibly change structure, not words ---
structure is what the rubric measures.
\end{solution}
\begin{exercise}[hard]\label{exr:ea9}
Argue for or against: ``partial credit write-ups should state what the
author could not prove.''
\end{exercise}
\begin{solution}
For: naming the gap is verifiable progress and distinguishes honest
incompleteness from bluffing; graders reward it. The cost is honesty
about scope.
\end{solution}
\begin{exercise}[hard]\label{exr:ea10}
Simulate a 4.5-hour, 6-proof paper: write the budget table, run two
problems under it, and report where the plan broke.
\end{exercise}
\begin{solution}
Most plans break at the deep-work cap --- the discipline being tested
is stopping, not solving.
\end{solution}
\begin{exercise}[easy]\label{exr:ea11}
A 3-problem exam is worth 21 points total (7 each). You solve $P_1$ in
25 min, $P_2$ in 50 min, and spend the last 15 checking. Sketch the
timeline.
\end{exercise}
\begin{solution}
$0$--$10$ triage/read; $10$--$35$ $P_1$; $35$--$85$ $P_2$;
$85$--$135$ $P_3$ attempt; $135$--$150$ checks --- the check block is
scheduled, never squeezed.
\end{solution}
\begin{exercise}[medium]\label{exr:ea12}
Rewrite this scratch into a submittable claim--lemma--proof: ``mod 7
squares are only 0,1,2,4, so $x^2+y^2=7z$ only if 7|x,y.''
\end{exercise}
\begin{solution}
\textbf{Claim.} $x^2+y^2\equiv0\pmod7 \Rightarrow 7\mid x,y$.
\textbf{Lemma.} $a^2\bmod7\in\{0,1,2,4\}$ (check $a=0..6$).
\textbf{Proof.} Pairs summing to $0\bmod7$ among
$\{0,1,2,4\}$: only $0+0$; hence $x^2\equiv y^2\equiv0$, and $7$
prime gives $7\mid x$, $7\mid y$.\ \qed
\end{solution}
\begin{exercise}[medium]\label{exr:ea13}
Two students claim ``I basically solved it'' after finding the right
method but writing no proof. Using the rubric, score each: (a) method
named, zero steps written; (b) method plus lemma statements, no
proofs.
\end{exercise}
\begin{solution}
(a) $0$---$1$: method recognition earns at most the setup point;
(b) $2$---$3$: correctly stated lemmas earn partial credit, but the
claim is unproven --- ``basically solved'' is a partial-credit
self-report, not a solve.
\end{solution}
\begin{exercise}[hard]\label{exr:ea14}
Design a verification protocol for a ``find all $f$'' solution that
found $\{f=0, f=\mathrm{id}\}$: list the three checks that catch the
two classic errors.
\end{exercise}
\begin{solution}
(1) \emph{Existence}: verify each member satisfies the equation
literally; (2) \emph{Completeness}: re-derive the family and confirm
no branch was dropped (e.g.\ $c^2=c$ allows both roots);
(3) \emph{Domain}: check the found functions are defined on all of
the stated domain. The classic errors are dropped branches and
unverified members.
\end{solution}
\begin{exercise}[hard]\label{exr:ea15}
Autopsy this fake-solve: ``Clearly $f$ is linear since additive
functions are lines, so $f(x)=kx$.'' Identify the two dishonest steps
and write the honest version.
\end{exercise}
\begin{solution}
Dishonest: ``clearly'' hides the hard part ($\mathbb R$ vs
$\mathbb Q$ domains), and ``additive functions are lines'' skips the
regularity hypothesis entirely. Honest: state the domain, derive
$f(q)=qf(1)$ on $\mathbb Q$, then name the regularity (continuity,
monotonicity, or boundedness) that extends it to $\mathbb R$.
\end{solution}

\section*{Notes}
Grading, fake-solve and checking conventions follow Chen's
\emph{Remarks on English}~\cite{chen-english} and the OTIS walkthrough
format~\cite{otis}; write-up mechanics follow Hammack~\cite{hammack}
and Velleman's technique summary~\cite{velleman}. Case-study problems
already cited: 2024 AIME I P10; 2021 AIME I P10; 2024 USAJMO P1;
USAMO 2003 rubric (olympiad.html)~\cite{chen-english}.

\section*{Further reading}
\begin{itemize}
\item E.~Chen, \emph{Remarks on English} --- how graders actually
read proofs; the reference for the write-up and fake-solve sections:
\url{https://web.evanchen.cc/handouts/english/english.pdf}
\item L.~Lamport, \emph{How to Write a 21st Century Proof} (paper)
--- the hierarchical proof structure the claim--lemma--proof style
descends from:
\url{https://lamport.azurewebsites.net/pubs/proof.pdf}
\item R.~Hammack, \emph{Book of Proof} --- the write-up
mechanics chapter:

\url{https://richardhammack.github.io/BookOfProof/}
\item D.~Velleman, \emph{How to Prove It} --- technique summaries for
each proof pattern the portfolio drills.
\item MIT OpenCourseWare, \emph{6.042J} ---
\url{https://ocw.mit.edu/6-042JF10}
--- the proof-writing recitations.
\item The OTIS walkthrough archive (E.~Chen's olympiad training
materials) ---
\url{https://web.evanchen.cc/otis.html}
--- model write-ups at contest
standard.
\end{itemize}

\begin{rulebox}[Watch instead]
\begin{itemize}
\item \textbf{Art of Problem Solving --- ``Work Backwards'' (Richard
Rusczyk)}, 12:00,
\url{https://www.youtube.com/watch?v=ZlKUhwNniOo}
--- the work-backwards
tactic that saves stalled exam problems. (Verified 2026-10-11.)
\item \textbf{Leslie Lamport --- ``How to Write a 21st Century
Proof''} (Heidelberg Laureate Forum 2017), $\sim$28 min,
\url{https://av.tib.eu/media/40137}
--- why claim--lemma--proof
structure beats prose proofs: the Write-up-standards section's
ancestor. (Verified 2026-10-11.)
\end{itemize}
\end{rulebox}

\chapternotes
